% =========================
% Chapter 6
% =========================
\chapter{Implementation}
\label{ch:implementation}

\section{Technology Stack}

\begin{table}[H]
\centering
\caption{Technology choices and their rationale.}
\label{tab:stack}
\begin{tabular}{p{3.2cm} p{3.6cm} p{7.2cm}}
\toprule
\textbf{Layer} & \textbf{Choice} & \textbf{Why} \\
\midrule
Server & FastAPI + uvicorn & Async WebSockets, single process serves API, WS, and frontend same-origin. \\
STT & faster-whisper (base) & Whisper weights on a CTranslate2 backend: ${\sim}4\times$ faster on CPU, no heavyweight framework dependency. \\
Audio decode & PyAV & Bundles FFmpeg in the wheel --- no system dependency for decoding PTT audio. \\
Detection (objects) & YOLOv8n & ${\sim}6$\,MB, CPU-real-time at 5\,FPS, COCO classes match the task vocabulary. \\
Detection (navigation) & YOLOv8m + 2 custom door models & Recall on indicators at angle/distance is the bottleneck; door models fill COCO's missing class. \\
Floor segmentation & SegFormer-B0 (ADE20K) & Real-time per-pixel walkability for the obstacle lane; optional --- degrades to named objects only. \\
Hand pose & MediaPipe Hands & 21 landmarks, CPU real-time, models bundled. \\
TTS & gTTS + LRU cache & Natural voice; caching hides latency for the fixed guidance phrases. \\
Command parsing & rapidfuzz & Structural prefix+noun fuzzy matching tolerates Whisper mishearings. \\
Frontend & Vanilla HTML/JS & No build step; one less moving part; served by the backend itself. \\
\bottomrule
\end{tabular}
\end{table}

\section{Backend Organization}

The backend is a single Python process organized in four layers.
\textbf{API layer} (\texttt{api/}): the WebSocket router dispatches JSON control
messages and tagged binary frames; the frame handler decodes JPEGs to arrays; the
audio handler runs the PTT pipeline (decode $\to$ transcribe $\to$ parse $\to$ FSM
event $\to$ spoken confirmation); the session object owns per-connection state.
\textbf{FSM layer} (\texttt{fsm/}): the five-state task machine with explicit
transitions and subscriber hooks. \textbf{Services layer} (\texttt{services/}):
speech, parsing, detection, hand pose, spatial reasoning, and guidance generation as
stateless, individually testable modules. \textbf{Task engines}: Object Allocation
(with Reach Guidance) and the Navigation package
(\texttt{services/navigation/}), each an asynchronous per-session loop started and
cancelled by FSM entry/exit hooks.

\section{The Per-Session Navigation Engine}
\label{sec:nav_engine}

The navigation package separates concerns strictly: \texttt{config.py} holds every
tuning constant in one place, each annotated with the failure that motivated its
value; \texttt{goals.py} the destination model; \texttt{perception.py} the door
funnel; \texttt{geometry.py} the compass and bearing mathematics;
\texttt{obstacles.py} the watchdog; \texttt{controller.py} the pure decision layer;
\texttt{models.py} lazy, thread-safe model loading shared across sessions; and
\texttt{engine.py} the asynchronous loop wiring them together. All mutable state
lives in a per-session \texttt{NavState} object, so concurrent users cannot
interfere --- a deliberate upgrade from the single-user prototype in which this
logic was first field-tested.

Three timing decisions matter for correctness. First, the engine measures each
tick's real duration and passes it to the controller, whose pacing constants are
expressed in \emph{seconds of observed frames} --- so timers keep their tuned
meaning regardless of load, and evidence thresholds (which are counts of
consecutive observations) remain counts. Second, the tick duration is clamped:
time during a stall contains no frames and therefore is not evidence that anything
appeared or disappeared. Third, a frozen-feed guard skips any frame that is stale
or already processed --- a locked screen or dropped connection can never accumulate
``evidence'' from a frozen image.

\section{The Speech Pipeline}
\label{sec:speech_pipeline}

Guidance text reaches the user through a per-session speech queue drained by a
dedicated worker: the perception loop enqueues and moves on, so TTS synthesis
latency never stalls detection or distorts its timers. The two-tier policy from
Section~\ref{sec:speech_discipline} is enforced at the queue: priority lines always
enqueue; ambient lines are dropped unless the pipe is idle. Synthesized MP3s are
cached in a bounded LRU keyed by exact text; at server startup an optional preload
warms both the navigation models and the fixed guidance phrases, so the first
navigation command of a session answers immediately and the common utterances
survive network hiccups. On the client, clips play through a single persistent,
gesture-primed \texttt{<audio>} element --- the pattern required for autoplaying
synthesized speech on iOS Safari.

\section{Frontend}
\label{sec:frontend}

The client is deliberately thin: five small JavaScript modules coordinate the
WebSocket, media capture (640$\times$480 JPEG at 5\,FPS; MediaRecorder for PTT),
the audio queue, a wake lock, a device-motion classifier (still / moving /
walking), and a compass monitor (throttled headings for the navigation scan). iOS
imposes the notable constraints: motion and orientation permissions must be
requested inside the Start button's click gesture, and audio must be unlocked by
playing a primed silent clip in that same gesture. The entire session UI is
voice-first --- the screen shows connection state, FSM state, and the last
transcription, but no interaction beyond Start, Stop, and the tap-anywhere
push-to-talk surface is required.

\section{Deployment}
\label{sec:deployment}

The backend serves the frontend same-origin, so a single HTTPS endpoint exposes the
whole system. For development and demos, a Cloudflare quick tunnel provides a real
HTTPS URL in one command --- necessary because browsers permit camera and
microphone access only over HTTPS (or localhost):

\begin{lstlisting}
cd backend && uvicorn main:app --port 8000
cloudflared tunnel --url http://localhost:8000
\end{lstlisting}

The phone opens the printed URL; no certificates, DNS, or router configuration are
involved. The heavyweight models load lazily on first use (or eagerly via the
preload flag); the custom door detector and verifier weights are committed to the
repository so a fresh clone runs the complete system.

\section{Test Infrastructure}
\label{sec:test_infrastructure}

All decision logic is written as pure, clock-injected code: trackers and controllers
receive detections, timestamps, and tick durations as arguments and return actions
and phrases --- no I/O, no models, no global clock. This is what makes the 286-test
suite possible: FSM transitions, command parsing across 50+ realistic transcriptions
(including mishearings), spatial bucketing with per-class priors, reach-guidance
state transitions, the Object Allocation tracker end-to-end with scripted frames and
a fake clock, the navigation goal model and arrival rules, the compass/bearing
mathematics (including the $\pm180^\circ$ seam), and complete scripted navigation
journeys --- scan, door choice, approach, obstacle preemption, doorway transit, new
room, semantic arrival --- run entirely from faked perception. The suite avoids
loading any ML model and completes in under one second, so it runs on every change.
Heavy perception (the door funnel against real pixels) is validated by the live
trials of Chapter~\ref{ch:evaluation}.
